\begin{resumo}

Este trabalho propõe a avaliação comparativa de modelos multimodais aplicados à extração estruturada de informações em provas objetivas manuscritas. Neste trabalho, o termo provas objetivas manuscritas refere-se a avaliações em papel cuja estrutura objetiva é preservada, mas cujos elementos visuais podem estar escritos manualmente, incluindo identificação do aluno, enunciados, numeração das questões, alternativas e marcações de resposta. O foco está na área de Sistemas de Informação, considerando a construção de um fluxo padronizado para captura, envio, processamento, validação e análise das respostas produzidas por modelos multimodais. O aplicativo móvel não executa modelos localmente; sua função é capturar imagens e apresentar resultados. O processamento ocorre no backend, que envia as mesmas imagens para um modelo multimodal aberto e um modelo multimodal proprietário ambos acessados via API. A metodologia adota Design Science Research para orientar o desenvolvimento do artefato e um experimento comparativo para avaliar acurácia, validade da saída estruturada, latência, custo operacional, robustez e necessidade de intervenção humana. Espera-se que o resultado indique qual abordagem técnica é mais adequada ao cenário estudado, considerando o desempenho dos modelos, a confiabilidade da saída JSON e os trade-offs envolvidos na integração de modelos multimodais a sistemas de informação.

\vspace{\onelineskip}

\noindent
\textbf{Palavras-chave}: Modelos multimodais. Sistemas de Informação. Processamento de documentos. JSON. Validação humana.

\end{resumo}
